\documentclass[UTF8]{article}
\usepackage{layout}
\hypersetup{pdftitle={2026全球开源发展报告·精要版｜第六章 v0.4},pdfsubject={十二项关键洞察、十大前瞻与十大战略建议}}
\begin{document}
\begin{canvas}
\Header{第六章・洞察、前瞻与战略建议}{blue}\Footer{01}
\pdfbookmark[1]{6.1 A 全球协作与体系能力}{p1}
\Box{15}{29}{180}{\fontsize{20}{26}\selectfont\color{navy}\sffamily \bfseries 6.1　十二项关键洞察}
\Box{15}{43}{151}{\fontsize{14}{18}\selectfont\color{blue}\sffamily \bfseries A　全球协作与体系能力}
\Box{168}{44}{27}{\fontsize{9.3}{13}\selectfont\color{muted}\sffamily 洞察 01—03}
\Box{15}{54}{180}{\fontsize{10.15}{15.5}\selectfont\color{muted}人工智能加快创造，开放协作让成果被更多人使用和改进。本报告把人才、组织、项目、代码、模型与数据联系起来观察，重点回答：谁在参与，资源怎样配合，以及哪些共同工作需要长期支持。}
\node[anchor=north west,inner sep=0pt] at \Pos{15}{71}{\includegraphics[width=180mm]{.build/assets/6-1-global.pdf}};
\Box{15}{185}{180}{\fontsize{8.1}{11.3}\selectfont\color{muted}\sffamily 图 6-1　从参与到持续协作。①②③统计期为 2025.07—2026.06。①三项占比分别计算；②OpenRank 根据协作关系衡量影响，145 个项目中 143 个有记录；③仓库是资源发布单元；④2026.10.16 局部模型样本。来源：图 2-8、4-3、5-1、5-4。}
\Box{15}{201}{180}{\fontsize{11.1}{14.7}\selectfont\color{blue}\sffamily \bfseries 01　把分散贡献组织起来，才能形成持续协作}
\Box{15}{207}{180}{\fontsize{10.15}{14.7}\selectfont\color{ink}在固定 29 国样本中，中美占活跃人数的 34.1\%，占项目影响力的 89.1\%。这提示：扩大参与的同时，也要培育能够组织协作的项目和团队，让参与者知道可以做什么，并有人回应问题、评审修改和持续维护。}
\Box{15}{227}{180}{\fontsize{11.1}{14.7}\selectfont\color{blue}\sffamily \bfseries 02　人工智能的发展，需要整套技术共同支撑}
\Box{15}{233}{180}{\fontsize{10.15}{14.7}\selectfont\color{ink}模型要与数据、软件工具和计算设备配合，才能进入实际应用。技术全景中，支撑模型训练和运行的基础设施占影响分值的 47.7\%。发展新应用，也需要长期建设这些共同使用的工具和设施。}
\Box{15}{253}{180}{\fontsize{11.1}{14.7}\selectfont\color{blue}\sffamily \bfseries 03　改进已有成果，也能创造新的价值}
\Box{15}{259}{180}{\fontsize{10.15}{14.7}\selectfont\color{ink}在局部模型网络中，552 个版本是在已有模型上改进而来，后来又被其他版本使用。让模型更适合具体任务、整理可用数据、做好测试和维护，都能为后续创新提供帮助，也是值得重视的开源贡献。}
\end{canvas}
\begin{canvas}
\Header{第六章・洞察、前瞻与战略建议}{china}\Footer{02}
\pdfbookmark[1]{6.1 B 中国模型开源与全球协作}{p2}
\Box{15}{29}{180}{\fontsize{20}{26}\selectfont\color{navy}\sffamily \bfseries 6.1　十二项关键洞察}
\Box{15}{43}{151}{\fontsize{14}{18}\selectfont\color{china}\sffamily \bfseries B　中国模型开源与全球协作}
\Box{168}{44}{27}{\fontsize{9.3}{13}\selectfont\color{muted}\sffamily 洞察 04—06}
\Box{15}{54}{180}{\fontsize{10.15}{15.5}\selectfont\color{muted}中国的开源贡献包括模型，也包括训练、运行和改进模型所需的软件工具。把中国发布的成果与全球使用情况联系起来，才能看清这些资源如何帮助其他团队继续创新。}
\node[anchor=north west,inner sep=0pt] at \Pos{15}{71}{\includegraphics[width=180mm]{.build/assets/6-2-china.pdf}};
\Box{15}{185}{180}{\fontsize{8.1}{11.3}\selectfont\color{muted}\sffamily 图 6-2　中国成果与全球创新。①2026.06 模型网络影响榜，按版本计席位；②2026.09.13 可识别跨国关系，覆盖有限；③CANN 统计期为 2025.07—2026.06，已剔除机器人；④沿用第三章项目实例，只说明功能。来源：图 3-14、3-17、3-20、5-5、5-6。}
\Box{15}{201}{180}{\fontsize{11.1}{14.7}\selectfont\color{china}\sffamily \bfseries 04　中国模型成为全球创新的重要起点}
\Box{15}{207}{180}{\fontsize{10.15}{14.7}\selectfont\color{ink}在可识别的跨国关系中，中国模型被美国发布者用于后续版本的关系占 42.4\%，是最大的单一方向。Qwen 的基础模型、DeepSeek 的推理示例等资源被继续使用，说明中国贡献还体现在帮助全球团队开展后续创新。}
\Box{15}{227}{180}{\fontsize{11.1}{14.7}\selectfont\color{china}\sffamily \bfseries 05　中国开源已覆盖模型、工具和计算软件}
\Box{15}{233}{180}{\fontsize{10.15}{14.7}\selectfont\color{ink}CANN 提供昇腾计算软件，verl 支持模型训练，torch\_npu 帮助 PyTorch 模型开发工具使用昇腾设备。中国已在多个技术环节提供开源成果。加强这些环节的配合，并持续更新和维护，有助于形成更完整的技术能力。}
\Box{15}{253}{180}{\fontsize{11.1}{14.7}\selectfont\color{china}\sffamily \bfseries 06　清晰的规则，让跨国合作更容易持续}
\Box{15}{259}{180}{\fontsize{10.15}{14.7}\selectfont\color{ink}合作伙伴需要知道资源能否使用和修改、材料是否齐全，以及遇到问题由谁负责。中国的开源制度建设可以从这些问题入手，完善授权、评价和共同维护规则，让国内外伙伴更容易理解条件、参与合作。}
\end{canvas}
\begin{canvas}
\Header{第六章・洞察、前瞻与战略建议}{violet}\Footer{03}
\pdfbookmark[1]{6.1 C 全球南方与可持续公共能力}{p3}
\Box{15}{29}{180}{\fontsize{20}{26}\selectfont\color{navy}\sffamily \bfseries 6.1　十二项关键洞察}
\Box{15}{43}{151}{\fontsize{14}{18}\selectfont\color{violet}\sffamily \bfseries C　全球南方与可持续公共能力}
\Box{168}{44}{27}{\fontsize{9.3}{13}\selectfont\color{muted}\sffamily 洞察 07—09}
\Box{15}{54}{180}{\fontsize{10.15}{15.5}\selectfont\color{muted}AI 普惠，要让更多地方用得起、用得好，并能自主改进和长期使用人工智能。全球南方主要指亚洲、非洲、拉美等地的发展中经济体。各地语言和需求不同，可以选择适合自己的创新方向。}
\node[anchor=north west,inner sep=0pt] at \Pos{15}{71}{\includegraphics[width=180mm]{.build/assets/6-3-public.pdf}};
\Box{15}{185}{180}{\fontsize{8.1}{11.3}\selectfont\color{muted}\sffamily 图 6-3　参与创新与持续使用的条件。①2026.09.13 跨国关系样本，选列三国；②③OpenGauge 889 条模型记录，材料较完整为 L3—L5、权限较开放为 O2—O4，沿用来源分级；④PyTorch 官方机制。②③④于 2026.11.11 期后核对。来源：图 5-5、5-9、4-13。}
\Box{15}{201}{180}{\fontsize{11.1}{14.7}\selectfont\color{violet}\sffamily \bfseries 07　全球南方也能成为创新成果的提供者}
\Box{15}{207}{180}{\fontsize{10.15}{14.7}\selectfont\color{ink}样本中，印度和巴西较多基于已有模型发布新版本，印度尼西亚也提供被他人继续使用的模型。各地可以从本地语言、数据和行业需求出发，逐步形成自己的成果，并为全球创新作出贡献。}
\Box{15}{227}{180}{\fontsize{11.1}{14.7}\selectfont\color{violet}\sffamily \bfseries 08　允许使用，还要让人能够独立使用和改进}
\Box{15}{233}{180}{\fontsize{10.15}{14.7}\selectfont\color{ink}617 条权限较开放的记录中，只有 117 条同时具备较完整的材料，占 19.0\%。使用者除了获得许可，还需要知道怎样训练、运行和改进模型。配套材料、可承担的成本、本地适用性与长期维护，都是普惠的必要条件。}
\Box{15}{253}{180}{\fontsize{11.1}{14.7}\selectfont\color{violet}\sffamily \bfseries 09　基金会有助于支持长期、共同需要的服务}
\Box{15}{259}{180}{\fontsize{10.15}{14.7}\selectfont\color{ink}PyTorch 的云资源计划将资金和计算资源用于共同测试，并安排技术支持。基金会通过这种方式汇集企业、研究者与社区的力量，支持大家都需要的基础工作。普惠成效还要看这些服务能否被更多地方持续使用。}
\end{canvas}
\begin{canvas}
\Header{第六章・洞察、前瞻与战略建议}{teal}\Footer{04}
\pdfbookmark[1]{6.1 D 中国产业、人才与城市结构}{p4}
\Box{15}{29}{180}{\fontsize{20}{26}\selectfont\color{navy}\sffamily \bfseries 6.1　十二项关键洞察}
\Box{15}{43}{151}{\fontsize{14}{18}\selectfont\color{teal}\sffamily \bfseries D　中国产业、人才与城市结构}
\Box{168}{44}{27}{\fontsize{9.3}{13}\selectfont\color{muted}\sffamily 洞察 10—12}
\Box{15}{54}{180}{\fontsize{10.15}{15.5}\selectfont\color{muted}中国不同团队提供不同资源，不同城市也形成了各自的优势。把产业分工、人才增长和地域分布放在一起看，有助于找到合作伙伴，也有助于地方选择适合自身条件的发展方向。}
\node[anchor=north west,inner sep=0pt] at \Pos{15}{71}{\includegraphics[width=180mm]{.build/assets/6-4-cities.pdf}};
\Box{15}{185}{180}{\fontsize{8.1}{11.3}\selectfont\color{muted}\sffamily 图 6-4　中国资源分工与城市互补。①③为 2026.10.17 组织样本累计下载，按发布方归属统计，不代表用户所在地；②固定十城人才样本选列五城，非全国完整排名，筛选总体仍待核。省域数据不能直接代表单个城市。来源：图 2-12、5-2、5-3。}
\Box{15}{201}{180}{\fontsize{11.1}{14.7}\selectfont\color{teal}\sffamily \bfseries 10　模型之外，同样存在重要的产业机会}
\Box{15}{207}{180}{\fontsize{10.15}{14.7}\selectfont\color{ink}中国样本中，模型与数据集下载排名前列的团队并不相同。基础模型、数据整理、软件工具和测试服务，各有专业分工。同时关注这些团队，才能发现模型排行榜之外的合作机会与产业发展空间。}
\Box{15}{227}{180}{\fontsize{11.1}{14.7}\selectfont\color{teal}\sffamily \bfseries 11　人才增长还要转化为持续解决问题的能力}
\Box{15}{233}{180}{\fontsize{10.15}{14.7}\selectfont\color{ink}固定十城中，北京与上海占人才样本的 52.5\%；杭州最近一年增长 17.4\%，上海净增 118 人。扩大人才队伍的同时，也需要培养愿意长期评审、测试、改进和维护项目的人，让人才增长带来持续的项目成果。}
\Box{15}{253}{180}{\fontsize{11.1}{14.7}\selectfont\color{teal}\sffamily \bfseries 12　城市可以通过专业分工和跨地区合作共同发展}
\Box{15}{259}{180}{\fontsize{10.15}{14.7}\selectfont\color{ink}浙江占中国样本模型下载的 56.8\%，北京与上海占数据下载的 75.4\%。不同资源集中在不同地方，为跨地区合作提供了线索。地方可以发挥自身所长，再与其他地区的模型、数据和应用团队合作，补充所需能力。}
\end{canvas}
\begin{canvas}
\Header{第六章・洞察、前瞻与战略建议}{blue}\Footer{05}
\pdfbookmark[1]{6.2　十大前瞻}{p5}
\Box{15}{29}{180}{\fontsize{20}{26}\selectfont\color{navy}\sffamily \bfseries 6.2　十大前瞻}
\Box{15}{43}{180}{\fontsize{12.4}{17}\selectfont\color{blue}\sffamily \bfseries 未来值得关注的十个变化}
\Box{15}{54}{180}{\fontsize{10.5}{16}\selectfont\color{muted}根据前五章的观察，提出近期、中期和面向 2030 的十项判断。这些前瞻说明可能的发展方向，能否实现仍取决于技术进步与合作条件。}
\Rect{15}{67}{180}{8}{paper}
\Box{18}{68.5}{174}{\fontsize{9.3}{13}\selectfont\color{muted}\sffamily 近期：发布后 6—12 个月　｜　中期：1—3 年　｜　2030：看长期方向与实现条件}
\Box{15}{79.5}{11}{\fontsize{18}{21}\selectfont\color{blue}\sffamily \bfseries 01}
\Box{15}{88.0}{11}{\fontsize{7.8}{10.5}\selectfont\color{blue}\sffamily 近期}
\Box{28}{80.0}{73}{\fontsize{11.9}{16.5}\selectfont\color{navy}\sffamily \bfseries 数据与应用将成为开源增长重点}
\Box{15}{94.0}{86}{\fontsize{10.5}{15.8}\selectfont\color{ink}报告统计期内，数据集和应用空间的新增增速高于模型。今后，更多团队有望围绕可用数据、软件工具和具体应用提供成果，丰富模型之外的开放资源。}
\Rule{15}{116.8}{101}{line}
\Box{109}{79.5}{11}{\fontsize{18}{21}\selectfont\color{blue}\sffamily \bfseries 02}
\Box{109}{88.0}{11}{\fontsize{7.8}{10.5}\selectfont\color{blue}\sffamily 近期}
\Box{122}{80.0}{73}{\fontsize{11.9}{16.5}\selectfont\color{navy}\sffamily \bfseries 创造越快，评审与维护越重要}
\Box{109}{94.0}{86}{\fontsize{10.5}{15.8}\selectfont\color{ink}AI 生成的成果增多后，项目需要更好地检查质量、处理问题和持续更新。能够组织这些工作的团队，有望把快速创造的成果转化为长期可用的资源。}
\Rule{109}{116.8}{195}{line}
\Box{15}{118.7}{11}{\fontsize{18}{21}\selectfont\color{blue}\sffamily \bfseries 03}
\Box{15}{127.2}{11}{\fontsize{7.8}{10.5}\selectfont\color{blue}\sffamily 近期}
\Box{28}{119.2}{73}{\fontsize{11.9}{16.5}\selectfont\color{navy}\sffamily \bfseries 中国模型有望支持更多全球应用}
\Box{15}{133.2}{86}{\fontsize{10.5}{15.8}\selectfont\color{ink}中国模型已成为许多后续版本的起点。若开放合作持续，更多海外团队有望将其用于本地语言和专业任务，中国贡献也将通过这些应用与改进继续扩大。}
\Rule{15}{156.0}{101}{line}
\Box{109}{118.7}{11}{\fontsize{18}{21}\selectfont\color{blue}\sffamily \bfseries 04}
\Box{109}{127.2}{11}{\fontsize{7.8}{10.5}\selectfont\color{blue}\sffamily 近期}
\Box{122}{119.2}{73}{\fontsize{11.9}{16.5}\selectfont\color{navy}\sffamily \bfseries 让模型适用不同设备将更重要}
\Box{109}{133.2}{86}{\fontsize{10.5}{15.8}\selectfont\color{ink}把模型用于不同设备和实际业务，需要解决软件兼容和运行效率等问题。帮助模型、工具与设备配合的团队，有望获得更多合作机会。}
\Rule{109}{156.0}{195}{line}
\Box{15}{157.9}{11}{\fontsize{18}{21}\selectfont\color{blue}\sffamily \bfseries 05}
\Box{15}{166.4}{11}{\fontsize{7.8}{10.5}\selectfont\color{blue}\sffamily 近期}
\Box{28}{158.4}{73}{\fontsize{11.9}{16.5}\selectfont\color{navy}\sffamily \bfseries 数据和配套材料的价值将更突出}
\Box{15}{172.4}{86}{\fontsize{10.5}{15.8}\selectfont\color{ink}模型越丰富，使用者越需要可靠数据、清楚的训练方法与可运行的软件环境。能够提供完整材料和使用指导的团队，将更有条件推动成果进入实际应用。}
\Rule{15}{195.2}{101}{line}
\Box{109}{157.9}{11}{\fontsize{18}{21}\selectfont\color{blue}\sffamily \bfseries 06}
\Box{109}{166.4}{11}{\fontsize{7.8}{10.5}\selectfont\color{blue}\sffamily 近期}
\Box{122}{158.4}{73}{\fontsize{11.9}{16.5}\selectfont\color{navy}\sffamily \bfseries 开放评价将更重视自主改进能力}
\Box{109}{172.4}{86}{\fontsize{10.5}{15.8}\selectfont\color{ink}评价开放程度时，除了能否下载和使用，还会更关注使用者能否理解、继续训练和独立维护。清晰授权与完整材料，有望成为长期合作的重要依据。}
\Rule{109}{195.2}{195}{line}
\Box{15}{197.10000000000002}{11}{\fontsize{18}{21}\selectfont\color{teal}\sffamily \bfseries 07}
\Box{15}{205.60000000000002}{11}{\fontsize{7.8}{10.5}\selectfont\color{teal}\sffamily 中期}
\Box{28}{197.60000000000002}{73}{\fontsize{11.9}{16.5}\selectfont\color{navy}\sffamily \bfseries 整套技术的协作将更影响竞争力}
\Box{15}{211.60000000000002}{86}{\fontsize{10.5}{15.8}\selectfont\color{ink}模型、数据、工具、计算设施与人才需要共同发挥作用。能够组织跨团队合作、让技术持续可用的地区和机构，有望形成更稳定的发展优势。}
\Rule{15}{234.40000000000003}{101}{line}
\Box{109}{197.10000000000002}{11}{\fontsize{18}{21}\selectfont\color{teal}\sffamily \bfseries 08}
\Box{109}{205.60000000000002}{11}{\fontsize{7.8}{10.5}\selectfont\color{teal}\sffamily 中期}
\Box{122}{197.60000000000002}{73}{\fontsize{11.9}{16.5}\selectfont\color{navy}\sffamily \bfseries 长期维护能力将更有价值}
\Box{109}{211.60000000000002}{86}{\fontsize{10.5}{15.8}\selectfont\color{ink}开源成果进入更多业务后，修复问题、保持兼容、培养接班人会更重要。能够长期承担这些工作的人才、团队和研究机构，将在生态中发挥更大作用。}
\Rule{109}{234.40000000000003}{195}{line}
\Box{15}{236.3}{11}{\fontsize{18}{21}\selectfont\color{teal}\sffamily \bfseries 09}
\Box{15}{244.8}{11}{\fontsize{7.8}{10.5}\selectfont\color{teal}\sffamily 中期}
\Box{28}{236.8}{73}{\fontsize{11.9}{16.5}\selectfont\color{navy}\sffamily \bfseries 全球南方将发展多样的本地创新}
\Box{15}{250.8}{86}{\fontsize{10.5}{15.8}\selectfont\color{ink}语言、教育、产业和公共服务需求，将推动各地选择自己的创新方向。本地团队若获得持续支持，有望在使用和改进已有成果的基础上，提供新的模型、数据和应用。}
\Box{109}{236.3}{11}{\fontsize{18}{21}\selectfont\color{violet}\sffamily \bfseries 10}
\Box{109}{244.8}{11}{\fontsize{7.8}{10.5}\selectfont\color{violet}\sffamily 2030}
\Box{122}{236.8}{73}{\fontsize{11.9}{16.5}\selectfont\color{navy}\sffamily \bfseries AI 普惠更依靠本地长期建设}
\Box{109}{250.8}{86}{\fontsize{10.5}{15.8}\selectfont\color{ink}面向 2030，普惠成效将更取决于本地团队能否自主选择、改进和维护技术。配套材料、人才与运行条件若得到长期支持，更多地区才有机会持续受益。}
\end{canvas}
\begin{canvas}
\Header{第六章・洞察、前瞻与战略建议}{teal}\Footer{06}
\pdfbookmark[1]{6.3　十大战略建议}{p6}
\Box{15}{29}{180}{\fontsize{20}{26}\selectfont\color{navy}\sffamily \bfseries 6.3　十大战略建议}
\Box{15}{43}{180}{\fontsize{12.4}{17}\selectfont\color{teal}\sffamily \bfseries 推动全球协作，让开放成果持续惠及更多人}
\Box{15}{54}{180}{\fontsize{10.5}{16}\selectfont\color{muted}围绕参与、供给、复用、治理与公共价值，提出面向政府、企业、研究机构和社区的十项建议，重点支持开放成果的共同建设、持续使用与长期维护。}
\Rule{15}{73}{195}{line}
\Box{15}{79.5}{11}{\fontsize{18}{21}\selectfont\color{teal}\sffamily \bfseries 01}
\Box{28}{80.0}{73}{\fontsize{11.9}{16.5}\selectfont\color{navy}\sffamily \bfseries 统筹人才、项目与资源的发展}
\Box{15}{90.0}{86}{\fontsize{11}{17}\selectfont\color{ink}制定开源战略时，把人才、组织、项目、代码、模型与数据放在一起看。识别已有优势，找到需要补充的环节，让政策和产业投入支持它们共同发展。}
\Rule{15}{116.8}{101}{line}
\Box{109}{79.5}{11}{\fontsize{18}{21}\selectfont\color{teal}\sffamily \bfseries 02}
\Box{122}{80.0}{73}{\fontsize{11.9}{16.5}\selectfont\color{navy}\sffamily \bfseries 为公共资源维护提供长期支持}
\Box{109}{90.0}{86}{\fontsize{11}{17}\selectfont\color{ink}把数据整理、文档、测试与维护纳入公共资助和企业投入，为承担这些工作的人提供稳定支持，让已经开放的成果能够持续使用、更新和改进。}
\Rule{109}{116.8}{195}{line}
\Box{15}{118.7}{11}{\fontsize{18}{21}\selectfont\color{teal}\sffamily \bfseries 03}
\Box{28}{119.2}{73}{\fontsize{11.9}{16.5}\selectfont\color{navy}\sffamily \bfseries 围绕共同需求深化全球合作}
\Box{15}{129.2}{86}{\fontsize{11}{17}\selectfont\color{ink}围绕已被跨国使用的模型、软件和数据，组织联合改进、测试与维护。通过共同解决问题扩大中国贡献、吸收全球经验，建立稳定的合作关系。}
\Rule{15}{156.0}{101}{line}
\Box{109}{118.7}{11}{\fontsize{18}{21}\selectfont\color{teal}\sffamily \bfseries 04}
\Box{122}{119.2}{73}{\fontsize{11.9}{16.5}\selectfont\color{navy}\sffamily \bfseries 补齐开放成果的配套材料}
\Box{109}{129.2}{86}{\fontsize{11}{17}\selectfont\color{ink}支持团队提供清楚的训练方法、数据说明和运行环境，邀请外部团队独立验证。帮助使用者理解成果、掌握方法，并在自己的条件下继续改进。}
\Rule{109}{156.0}{195}{line}
\Box{15}{157.9}{11}{\fontsize{18}{21}\selectfont\color{teal}\sffamily \bfseries 05}
\Box{28}{158.4}{73}{\fontsize{11.9}{16.5}\selectfont\color{navy}\sffamily \bfseries 把使用条件和维护责任说清楚}
\Box{15}{168.4}{86}{\fontsize{11}{17}\selectfont\color{ink}明确资源可以怎样使用和修改、哪些材料能够获得、由谁负责维护。评价结果要有依据，争议要能解释，错误要能纠正，让合作伙伴据此作出判断。}
\Rule{15}{195.2}{101}{line}
\Box{109}{157.9}{11}{\fontsize{18}{21}\selectfont\color{teal}\sffamily \bfseries 06}
\Box{122}{158.4}{73}{\fontsize{11.9}{16.5}\selectfont\color{navy}\sffamily \bfseries 共建大家都能使用的测试设施}
\Box{109}{168.4}{86}{\fontsize{11}{17}\selectfont\color{ink}由企业、基金会和高校共同提供测试环境，帮助模型和软件适用不同设备与业务场景。共享计算资源与技术支持，降低团队参与和使用的门槛。}
\Rule{109}{195.2}{195}{line}
\Box{15}{197.10000000000002}{11}{\fontsize{18}{21}\selectfont\color{teal}\sffamily \bfseries 07}
\Box{28}{197.60000000000002}{73}{\fontsize{11.9}{16.5}\selectfont\color{navy}\sffamily \bfseries 发挥地方优势，推动跨地区合作}
\Box{15}{207.60000000000002}{86}{\fontsize{11}{17}\selectfont\color{ink}地方根据人才、模型、数据和产业需求选择重点方向，支持跨城市共同项目。通过专业分工与资源共享，补充彼此需要的能力，减少重复建设。}
\Rule{15}{234.40000000000003}{101}{line}
\Box{109}{197.10000000000002}{11}{\fontsize{18}{21}\selectfont\color{teal}\sffamily \bfseries 08}
\Box{122}{197.60000000000002}{73}{\fontsize{11.9}{16.5}\selectfont\color{navy}\sffamily \bfseries 支持全球南方自主开展本地创新}
\Box{109}{207.60000000000002}{86}{\fontsize{11}{17}\selectfont\color{ink}由当地机构和社区确定需求、语言和使用方式，国际合作为其提供团队培养、计算资源与维护支持，让当地伙伴能够自主选择技术并持续改进。}
\Rule{109}{234.40000000000003}{195}{line}
\Box{15}{236.3}{11}{\fontsize{18}{21}\selectfont\color{teal}\sffamily \bfseries 09}
\Box{28}{236.8}{73}{\fontsize{11.9}{16.5}\selectfont\color{navy}\sffamily \bfseries 让长期协作与维护得到认可}
\Box{15}{246.8}{86}{\fontsize{11}{17}\selectfont\color{ink}在人才培养、考核和资助中，重视测试、文档、软件适配与维护等贡献。围绕真实项目培养人才，为持续解决问题的人提供发展机会与稳定支持。}
\Box{109}{236.3}{11}{\fontsize{18}{21}\selectfont\color{teal}\sffamily \bfseries 10}
\Box{122}{236.8}{73}{\fontsize{11.9}{16.5}\selectfont\color{navy}\sffamily \bfseries 把普惠成效作为公共投入的重点}
\Box{109}{246.8}{86}{\fontsize{11}{17}\selectfont\color{ink}评价公共投入时，同时关注资源规模、使用成本、本地适用性、服务持续性和实际受益情况，推动更多人用得起、用得好，并能够长期使用人工智能。}
\end{canvas}
\end{document}
